\chapter{Discusión} \label{sec:Discusion}

\section{El papel determinante del dataplane}

El rendimiento, la eficiencia de recursos y las capacidades de seguridad de un clúster Kubernetes no son propiedades del clúster en sí, sino del plugin CNI que implementa su plano de datos. Las diferencias observadas entre los cuatro CNIs ---con hardware, versión de Kubernetes y topología idénticos en cada ciclo--- lo confirman con claridad.

La ventaja de Calico no es marginal. Su latencia promedio (12.54 ms) es casi la mitad que la de Flannel (21.53 ms), y su estabilidad ---medida por el jitter (7.31 ms frente a 7.76 ms de Flannel)--- es ligeramente superior, aunque en retransmisiones Flannel resulta muy superior (11774 frente a 68993 de Calico). Esta brecha de rendimiento en latencia se debe a que Calico optimiza el enrutamiento directo y el socket local en el kernel, mientras que Flannel realiza una encapsulación VXLAN simple pero estable. Este comportamiento es consistente con los trabajos de Zhang et al.~\cite{ref27} y con la fundamentación de Pfaff et al.~\cite{ref22}, que anticipaban ventajas para arquitecturas avanzadas.

El liderazgo de Calico y Cilium en rendimiento tiene, sin embargo, un costo concreto en consumo de RAM bajo estrés: Calico requiere 1672.83 MB y Cilium 1497.19 MB, superando en un 42\% y 27\% respectivamente a Flannel (1176.54 MB). Este compromiso no es un defecto de implementación, sino consecuencia directa de la complejidad funcional: el mantenimiento de mapas eBPF de Cilium y el ruteo BGP/VXLAN de Calico con tablas dinámicas de conexiones residen en memoria. El modelo MCDA captura este compromiso correctamente, lo que explica por qué Cilium y Calico no lideran en el perfil IoT, donde la eficiencia de recursos es crítica.

\section{Flannel: el riesgo silencioso de la configuración predeterminada}

El hallazgo más crítico en seguridad no es la diferencia de latencia entre CNIs, sino el comportamiento de Flannel ante las NetworkPolicies. Las políticas se crean en la API de Kubernetes sin ningún error, el sistema las acepta, pero el tráfico que debería bloquearse continúa fluyendo sin restricción. Un operador que despliegue Flannel y aplique políticas de aislamiento obtiene una falsa certeza de seguridad: el clúster parece configurado correctamente, pero en la práctica está expuesto.

El dato cobra especial peso porque Flannel es el CNI predeterminado de K3s y de numerosas guías de inicio rápido de Kubernetes, lo que lo convierte en la opción más frecuente en entornos de aprendizaje y en primeras implementaciones productivas. La literatura advertía sobre esta limitación~\cite{ref4,ref18}, pero la evidencia experimental confirma que el riesgo no es teórico: en todos los casos de prueba que requerían enforcement de políticas, Flannel los reprobó. Cualquier organización que opere Flannel en producción con NetworkPolicies definidas debe tratar esas políticas como declarativas ---sin efecto real sobre el tráfico--- hasta realizar una validación explícita del enforcement.

\section{Zero Trust es viable sin penalizar el rendimiento}

Una premisa frecuente en los equipos de operaciones es que activar políticas de red completas degrada el rendimiento del clúster de forma inaceptable. Los resultados contradicen esa premisa con datos concretos. El overhead de latencia promedio al activar políticas de seguridad fluctuó entre reducciones significativas de latencia (de hasta el $-43.16\%$ en Cilium Zero Trust) debido a optimizaciones de canal, y un incremento máximo del $+31.45\%$ en Calico Multi-Tier. En cuanto a throughput, el impacto negativo se mantuvo por debajo del $4.35\%$ en Calico y Antrea, mientras que Cilium y Antrea experimentaron mejoras debido al bypass del datapath en el kernel.

Este resultado contrasta favorablemente con las advertencias de Chandramouli et al.~\cite{ref26} sobre el costo de las arquitecturas de Service Mesh ---que emplean proxies L7 en espacio de usuario---: las NetworkPolicies implementadas en el plano de datos del CNI operan en el kernel y son órdenes de magnitud más eficientes que un proxy intermedio. El overhead documentado en la literatura de Service Mesh no aplica a las políticas nativas del CNI. Los datos no justifican diferir la adopción de Zero Trust por temor al impacto en el rendimiento, siempre que el CNI soporte NetworkPolicies en su plano de datos.

\section{Alcance y limitaciones del modelo MCDA}

El modelo MCDA produjo recomendaciones diferenciadas por perfil de uso a partir de criterios heterogéneos. Que ningún CNI domine en todas las dimensiones de forma simultánea valida la necesidad del enfoque multicriterio: una comparación basada únicamente en throughput habría recomendado Calico para todos los perfiles, lo cual sería técnicamente incorrecto para el perfil IoT, donde Flannel ofrece un balance superior por su bajo consumo de recursos.

El modelo tiene, sin embargo, limitaciones concretas. Los pesos asignados a cada criterio por perfil son juicios de diseño razonados, no valores calibrados empíricamente a partir de organizaciones reales; un trabajo futuro podría derivarlos mediante entrevistas estructuradas con equipos de ingeniería. La evaluación de NetworkPolicies L7 se implementó como variable binaria (soporta/no soporta), perdiendo matices sobre la facilidad de expresar reglas complejas o el overhead de distintos patrones de filtrado. Finalmente, el experimento se ejecutó sobre dos nodos Worker; los efectos de escala con decenas de nodos y miles de Pods ---donde el overhead del plano de control del CNI se vuelve más visible--- podrían modificar el ranking, especialmente en las dimensiones de RAM y CPU en reposo.

\section{Respuesta al problema de investigación}

La pregunta de investigación era cómo comparar complementos de red y automatizar políticas de comunicación segura para optimizar recursos sin degradar el rendimiento. Los resultados ofrecen tres respuestas concretas.

La comparación objetiva es posible mediante un testbed reproducible, benchmarks automatizados y un modelo MCDA que pondera los criterios según el perfil de uso. El prototipo construido demuestra que esta comparación puede exponerse de forma interactiva y actualizable, sin exigir al usuario la interpretación de datos brutos ni conocimientos estadísticos.

La automatización de políticas de seguridad es viable sin sacrificar rendimiento: Calico, Cilium y Antrea lo hacen con un impacto de overhead controlado y moderado. La integración con GitOps garantiza la persistencia de estas políticas ante modificaciones no autorizadas, convirtiendo la seguridad en un estado mantenido continuamente y no en una configuración puntual.

La reducción del sobredimensionamiento tiene respaldo cuantitativo: la diferencia de consumo entre Flannel (19.45\% CPU y 1176.54 MB RAM) y Calico (29.05\% CPU y 1672.83 MB RAM) implica que, a escala de 50 nodos, es posible liberar el equivalente a aproximadamente 4.8 núcleos de CPU y 24.8 GB de memoria RAM, disminuyendo costos de hardware y mejorando la densidad de contenedores en el clúster sin modificar la carga de trabajo ni el rendimiento de las aplicaciones.
